\chapter{植物幼苗细粒度分类}
\label{chap:plant}

\section{任务定义与数据}
课程赛题要求为每张幼苗图像指定五个物种之一。训练集 500 张图像，五类各
100 张。测试集 378 张，输入图像分辨率和背景条件变化较大。典型图像包含植物叶片、
土壤及容器边缘。Black-grass 和 Loose Silky-bent 均呈细长的草叶形态，类别差异
集中在叶宽、分蘖和细节结构，单靠叶片数量或背景颜色难以稳定判断。课程图像来自
植物幼苗基准所涉及的数据域。原基准覆盖十二个物种，本文的五类数量和划分来自
课程数据目录\citep{giselsson2017plants}。

设 $x_i$ 为原始图像，$y_i\in\{1,\ldots,5\}$ 为类别。预测器输出 $5$ 维
概率向量，提交文件按样例的图像 ID 顺序写入 \code{ID,Category}。训练集采用
分层五折，每折约 400 张图像拟合、100 张图像验证。折外预测（out-of-fold, OOF）
仅拼接各图像所在验证折的输出，准确率与宏平均 F1 定义为
\begin{align}
  \mathrm{Acc}&=\frac{1}{n}\sum_{i=1}^n\mathbf 1[\hat y_i=y_i],\\
  \mathrm{MacroF1}&=\frac{1}{5}\sum_{c=1}^{5}
          \frac{2\,TP_c}{2\,TP_c+FP_c+FN_c}.
\end{align}
$TP_c,FP_c,FN_c$ 按一对其余类别计算。赛题标为 F1，历史得分粒度与 $1/378$
相符。由于平台聚合实现未随提交文件保存，本文将 Kaggle 原始得分和本地准确率
分别呈现。

\section{方法与结构}

\subsection{前景裁剪及坐标顺序}
图~\ref{fig:plant-arch} 给出从原始图像到最终类别的处理流程。透明 PNG 先合成到
背景以消除透明边缘。程序在 HSV 空间从绿色像素估计前景包围盒，向外留边并补齐
到正方形。裁剪首先在产生包围盒的坐标系中执行，随后缩放至骨干的输入尺寸。
去背景配置在裁剪之外使用前景掩膜弱化土壤纹理。HSV 阈值具有可解释性，低饱和度
的暗叶、非绿色幼叶和强光反射会影响掩膜质量。

\archdiagram{plant_arch}{幼苗分类的前景处理、迁移学习和集成流程}{fig:plant-arch}

旧版本有一条路径在缩放后的图像上沿用原图像素坐标。若原图宽度 $W$ 已缩放为
$W'$，原坐标 $u$ 对应 $uW'/W$，直接使用 $u$ 会把包围盒移到错误位置并可能越界。
Pillow 的越界填充产生大块黑区，模型最终接收的叶片信息减少。以暗像素占比筛查
异常图像后，程序回到原图坐标裁剪，再执行缩放。实际修复需要同步重建缓存，
否则旧缓存仍会把错误输入送进训练器。\ifLocalReport 图~\ref{fig:plant-case} 展示坐标错用如何改变输入。\fi
\privatechart{plant_case.png}{植物幼苗裁剪缺陷诊断示例}{fig:plant-case}
{案例来自原项目的缺陷诊断图。图中内容为预处理输入对照，性能变化见表~\ref{tab:plant-pairs}。}

\subsection{骨干和分类头}
主线使用 ImageNet 或 ImageNet-22K 预训练的 ConvNeXt-Tiny、ConvNeXt-Small 与
EfficientNet-B0\citep{liu2022convnext,tan2019efficientnet}。ConvNeXt 的阶段式
降采样提供从叶缘到整体株型的多尺度表征。深度卷积增加空间感受范围，逐点变换负责
通道融合。EfficientNet-B0 提供不同的卷积与通道注意力归纳偏置，适合比较较低
分辨率下的轻量基线。程序通过 timm 初始化骨干，保留对应的全局池化和
归一化，再将分类输出维度置为五。输入是 $B\times3\times R\times R$，输出为
$B\times5$。不同 $R$ 和不同骨干一起改变了训练成本，跨骨干的成绩差异涉及
预训练权重、参数量与输入尺度。

\subsection{分类目标、伪标签和集成}
记 $\mathbf p=\operatorname{softmax}(\mathbf z)$。交叉熵使用标签平滑
$\varepsilon=0.1$。当批次启用 Mixup 或 CutMix 时，两个类别的平滑目标按混合
系数 $\lambda$ 加权，令
\[
 q_c=\lambda( (1-\varepsilon)\mathbf 1[c=y_a]+\varepsilon/5)
      +(1-\lambda)((1-\varepsilon)\mathbf 1[c=y_b]+\varepsilon/5),
 \qquad \mathcal L=-\sum_{c=1}^5q_c\log p_c .
\]
Mixup 在像素空间插值，CutMix 拼接图块并依据实际保留面积计算权重。其作用是
降低对 500 张标注图像的过拟合。高置信度伪标签来自本项目对未标注测试图的
预测，阈值设为 0.9。典型修复轮次记录了 184 张入选样本。该流程属于传导式
设置，结论应限定到当前固定测试池。

测试时对同一图像的两个尺度、四个 $90^\circ$ 旋转和两个镜像状态预测，共
16 个视图。先在每个模型内平均 Softmax 概率，再在折与配置之间融合。
设 $\bar p_c(x)$ 为融合概率，类别先验校正使用
$\arg\max_c \bar p_c(x)\hat\pi_c^{\,\tau}$。
训练先验近似均匀，程序用 OOF 混淆矩阵估计类条件预测分布 $A$，并求一个
非负最小二乘近似 $A^\top\hat{\boldsymbol{\pi}}\approx\boldsymbol q$，
其中 $\boldsymbol q$ 来自测试预测分布。先验方法依赖类条件分布在训练与
测试间近似稳定。混淆矩阵中的两类高度相似时，反演对估计误差更敏感。

\section{实验设计与训练记录}
基线使用 EfficientNet-B0 和 288 像素输入。ConvNeXt 路线依次考察增强、
伪标签、去背景和不同输入尺度。每个配置复用五折。数据增强包括随机裁剪缩放、
水平及垂直翻转、旋转与随机擦除，部分配置加入 RandAugment 和软标签增强。
拍摄姿态变化支持旋转增强。强扰动有时会丢掉细长叶片，因此按骨干和分辨率
复核其结果。

训练先冻结骨干，仅用 AdamW 更新分类头，随后解冻模型并以 OneCycleLR
余弦退火微调\citep{loshchilov2019adamw}。日志和模型保存路径表明该流程评估
中间轮次，但训练脚本写出末轮模型，并以该权重生成折外预测。表~\ref{tab:plant-cfg}
按修复后实验的 \code{summary.json} 和逐折日志列出实际设置。每个配置的五折均
有与计划相同数量的记录。以 Tiny 伪标签配置第 0 折为例，44 轮完成，末轮验证
准确率为 0.960，见图~\ref{fig:plant-curve}。中间峰值与末轮权重分别记账，
便于检查选择规则。

\begin{table}[H]
\centering\small
\caption{修复后代表性配置及训练记录}\label{tab:plant-cfg}
\begin{tabularx}{\linewidth}{@{}Yrrrrr@{}}
\toprule
配置 & 输入边长 & 批大小 & 头阶段 & 微调 & OOF Acc \\
\midrule
EfficientNet-B0 & 288 & 32 & 4 & 46 & 0.922 \\
ConvNeXt-Tiny & 320 & 32 & 4 & 40 & 0.932 \\
ConvNeXt-Tiny，伪标签 & 320 & 32 & 4 & 40 & 0.942 \\
ConvNeXt-Small，伪标签和去背景 & 320 & 16 & 4 & 40 & 0.950 \\
ConvNeXt-Tiny，伪标签和去背景 & 384 & 24 & 4 & 40 & 0.928 \\
\bottomrule
\end{tabularx}
\end{table}
\chart{plant_training.pdf}{ConvNeXt-Tiny 伪标签配置第 0 折的训练记录}{fig:plant-curve}
{虚线标记第 4 轮结束，随后解冻并微调。左图训练目标受批次混合增强影响，右图为确定性验证输入的准确率。}

\section{结果与误差分析}
\subsection{修复前后及有效对照}
修复后，同为 ConvNeXt-Small、320 像素和去背景的配置，折外准确率由
0.9240 升至 0.9500，即 500 张图像中多分对 13 张。Tiny 的修复后 320
像素基线与伪标签版分别为 0.9320 和 0.9420。其差异对应 5 张图像，
说明在当前划分上有收益，仍需其他划分验证稳定性。两个同版本对照见表
\ref{tab:plant-pairs}，图~\ref{fig:plant-data}展示类别数和各模型组的结果。

\begin{table}[H]
\centering\small
\caption{裁剪修复前后的折外准确率与榜单记录}\label{tab:plant-pairs}
\begin{tabularx}{\linewidth}{@{}Yrrl@{}}
\toprule
比较对象 & 修复前 & 修复后 & 口径 \\
\midrule
ConvNeXt-Tiny，320，伪标签 & 0.9280 & 0.9420 & OOF Acc \\
ConvNeXt-Small，320，去背景和伪标签 & 0.9240 & 0.9500 & OOF Acc \\
旧六配置与新四配置集成 & 0.9300 & 0.9600 & OOF Acc \\
Kaggle \code{v3} 与 \code{v16} & 0.9523 & 0.9735 & Public / Private \\
\bottomrule
\end{tabularx}
\end{table}
\chart{plant_data.pdf}{五类训练样本数与修复后单组 OOF 准确率}{fig:plant-data}
{右图每一组由同一配置的五个折模型产生折外预测。横轴放大至 0.90 至 0.963 以显示差异。}

集成折外记录由 0.9300 提至 0.9600，同时涉及成员变化和伪标签重建。线上由
0.9523 到 0.9735，提高 2.12 个百分点。修复机制通过图像诊断和同类模型
对照得到支持，集成差值还包含其余改动。修复后的两配置和四配置、施加或省略
先验的若干提交均为 0.9735。相同得分允许存在不同逐样本预测，不能据此断言
两种推理配方完全等效。最终引用 \code{v16} 的提交编号为 \code{56983147}。本地 v13 与 v16 的 CSV 字节一致。
配置级训练结果可以逐项查证。旧记录对 \code{v16} 的权重配比和先验文件
对应不够完整，复现附录提供可执行的训练和推理入口以及本地成品哈希。

\subsection{细粒度混淆与同源数据风险}
旧阶段的逐类 F1 显示，Black-grass 与 Loose Silky-bent 分别为 0.854 与
0.871。其余三类为 0.970 至 1.000。这组数字对应较早的折外汇总，说明早期
错误集中在形态相似的两类。仅用全局尺度调整难以保留叶尖与叶脉，384 像素
配置的 OOF 也低于表现最好的 320 像素 Small。两者骨干和训练配方存在差异，
因此结果限定为配置级比较。

从公开同源植物库增加标注数据的支线取得 0.9947。原记录报告 378 张测试图中
查到 41 张逐像素副本，且扩充训练数据包含这批同源图像。外部库的域匹配与
样本重叠共同影响榜单分数。该结果保留为数据审计事实，主结论采用
0.9735 的课程标注加未标注图像伪标签路线。

\section{本章小结}
本任务通过修复裁剪坐标并重建缓存改善了输入质量。修复后的 Small 单配置 OOF 准确率为 0.9500。最终提交的两侧榜单均为 0.9735。
相似禾本科类别的混淆以及少样本折外评估的波动仍限制对新数据的外推。
